\section{Quadratic relation spaces and intrinsic reconstruction of pencils}
\label{sec:natural-pencils}

The multiplication construction itself supplies a wider natural class
for the coefficient criterion.  In this section the number of quadratic
relations is allowed to vary independently of the embedding dimension.
For two relations the resulting theorem recovers every pencil of
quadrics, including singular pencils.  No extra block is added to the
multiplication map.

\subsection{Automatic coefficient extraction for quadratic multiplication}

Let \(\dim V=n\ge3\), \(W=\Sym^2V\), \(N=n(n+1)/2\), and
\(R\in\Gr(r,W)\), where \(1\le r<N-n\).  Put \(p=N-r>n\),
\(S=W/R\), and
\[
 A_R=\C\oplus V\oplus S,\qquad
 X_R=\Gr(n,V\oplus S).
\]
The only nonzero product in the maximal ideal is the quotient map
\(\gamma:\Sym^2V\to S\).  Write \(\mathcal K\) for the
tautological bundle on \(X_R\), and use the same definition
\eqref{eq:failure-intro} of \(\widehat D_R\), with \(A_R\) in
place of \(B_R\).  Thus \(\widehat D_R\) is the failure scheme
of generating \(A_R\) by a varying \(n\)-plane and the unit.
The notation in this section agrees with the web construction when
\((n,r)=(4,4)\).

Let \(\Lambda_{n,p}\) denote the partitions occurring in the
classical multiplicity-free decomposition
\begin{equation}\label{eq:general-exterior-decomposition}
 F(V):=\bigwedge^p\Sym^2V
       =\bigoplus_{\lambda\in\Lambda_{n,p}}\mathbb S_\lambda V.
\end{equation}
Every partition has size \(2p\); terms of length greater than \(n\)
are omitted.  Choose a nonzero volume form on \(S\) and let
\(\beta_R\in F(V)^*\) be its pullback by \(\bigwedge^p\gamma\).
Its component in \(\mathbb S_\lambda V^*\) is denoted by
\(\beta_{R,\lambda}\).  Only the lines of the nonzero components
will be used, so the volume choice is immaterial.

\begin{theorem}[Automatic intrinsic coefficient extraction]
\label{thm:automatic-quadratic-coefficients}
For every \(R\in\Gr(r,\Sym^2V)\) in the stated range, the
abstract scheme \(\widehat D_R\) determines
\[
       \bigl(\C\beta_{R,\lambda}\bigr)_{\lambda\in\Lambda_{n,p}}
\]
up to one common projective transformation of \(V\), with zero
components recorded as zero subspaces.  More precisely, the intrinsic
construction recovers the base \(B=\Gr(n,S)\), its degree-one
normal-cone bundle, and the residual ideal generated in degree \(2p\)
with coefficient module
\begin{equation}\label{eq:automatic-residual-module}
 \mathcal J_{2p}=
 \bigoplus_{\lambda\in\Lambda_{n,p}}
       \C\beta_{R,\lambda}\otimes\mathbb S_\lambda\mathcal U
 \ \subset\ \Sym^{2p}(V^*\otimes\mathcal U),
\end{equation}
where \(\mathcal U\) is the tautological rank-\(n\) bundle on
\(B\).  Consequently all geometric hypotheses of
Theorem~\ref{thm:structural-coefficient-reconstruction} hold for this
entire Grassmannian family, not just for smooth webs.
\end{theorem}
\begin{proof}
On a frame chart of \(\mathcal K\), write its inclusion as
\(\binom M{A_0}\), with \(M\) an \(n\times n\) matrix.  Since
\(A_R\) is cube-zero, columns of degree greater than two vanish.
After retaining the scalar, linear, and quadratic columns, multiplication
has the matrix
\begin{equation}\label{eq:general-multiplication-matrix}
 \begin{pmatrix}
  1&0&0\\
  0&M&0\\
  0&A_0&\gamma\Sym^2M
 \end{pmatrix}.
\end{equation}
A nonzero maximal minor must use the scalar column and all \(n\)
linear columns: these are the only columns meeting the first \(1+n\)
rows.  Expanding along those rows gives the identity of ideals
\begin{equation}\label{eq:general-fitting-factorization}
 I_{\widehat D_R}=(\det M)I_p(\gamma\Sym^2M).
\end{equation}
The same expansion works for every \(m\ge2\).  Where \(M\) is
invertible, the quadratic map is surjective.  Hence the reduction of
\(\widehat D_R\) is exactly the Schubert divisor where \(M\)
is singular, independent of \(R\) after identifying the vector
spaces \(S\).

Successive reduced singular loci of that divisor are the projection
rank loci.  Their last nonempty member is \(B=\Gr(n,S)\).
For clarity, the usual determinantal assertion here is local: at a
projection-rank-\(k\) point, an invertible \(k\)-minor splits
away, leaving a generic square \((n-k)\)-matrix and smooth
coordinates.  Every remaining entry can vary independently by
moving the kernel of the projection in \(V\) modulo its image.
Thus its reduced singular locus is the next rank locus.  Since
\(p>n\), the final base is nonempty, connected and projective,
with \(H^0(B,\OO_B)=\C\).

Near \(B\), a Grassmannian chart consists of a varying plane in
\(S\) and its graph \(T:\mathcal U\to V\).  Formula
\eqref{eq:general-fitting-factorization} is independent of the other
graph coordinates and homogeneous in \(T\).  It therefore gives
its actual normal-cone ideal, not merely an initial reduced support:
\begin{equation}\label{eq:general-normal-cone}
 C_B\widehat D_R\subset\Tot\Hom(\mathcal U,V),\qquad
 \mathcal I=(\det T)\mathcal J,\quad
 \mathcal J=I_p(\gamma\Sym^2T).
\end{equation}
There are no degree-one equations.  The ambient bundle in
\eqref{eq:general-normal-cone} is consequently the dual of the
intrinsically recovered degree-one conormal.  Its reduced ideal is
the determinant ideal \(\mathcal D\).  The determinant polynomial
is monic for a suitable monomial order, so cancellation gives
\((\mathcal I:\mathcal D)=\mathcal J\).  This is an intrinsic
ideal quotient; choosing a generator of the determinant line is not
part of the recovered datum.

The rank-one locus of the reduced cone determines the two ruling
parameter schemes by Lemma~\ref{lem:rank-one-fano-scheme}.  They are
\(\Pj(V)\times B\) and \(\Pj_B(\mathcal U^*)\).
Exactly the first is projectively trivial.  Indeed a Schubert line in
\(\Gr(n,S)\) restricts \(\mathcal U^*\) to
\(\OO(1)\oplus\OO^{n-1}\).  A trivial projective bundle on
that line would make this bundle a line-bundle twist of
\(\OO^n\); its determinant would then have degree divisible by
\(n\), rather than degree one.  Thus the trivial ruling is
intrinsically oriented.  Lemma~\ref{lem:oriented-segre-descent}
applies: an abstract isomorphism gives a morphism \(B\to\PGL(V)\),
constant because \(B\) is connected projective and \(\PGL(V)\)
is affine.  Local linear lifts and right changes of frame give
\(T\mapsto gTh^{-1}\) for this one fixed \([g]\).

It remains to prove \eqref{eq:automatic-residual-module}, rather
than assume it as a condition on the example.  The minors are the
image, for fixed \(\beta_R\), of the universal coefficient map
\begin{equation}\label{eq:general-coefficient-map}
 F(V)^*\otimes F(U)\longrightarrow
 \Sym^{2p}(V^*\otimes U),\qquad
 \beta\otimes\xi\longmapsto[T\mapsto\beta(F(T)\xi)].
\end{equation}
In characteristic zero, the polynomial functor \(F\) has the
multiplicity-free decomposition \eqref{eq:general-exterior-decomposition};
see \cite[Theorem~3.5]{ChengWang}.  The Cauchy decomposition of the
target contains each diagonal
\(\mathbb S_\lambda V^*\otimes\mathbb S_\lambda U\)
once, and no off-diagonal pair.  Equivariance kills every
off-diagonal pair.  On each diagonal pair the coefficient map is
nonzero: identify \(U\) with \(V\), evaluate at \(T=1\), and
choose \(\beta,\xi\) in that summand with
\(\beta(\xi)\ne0\).  Schur's lemma then identifies its image
with the specified Cauchy summand.  This is the actual matrix-coefficient
inclusion, not merely a dimension count.  Keeping \(\beta_R\)
fixed proves \eqref{eq:automatic-residual-module}.  All generators
have degree \(2p\), and at least one minor is nonzero since
\(\gamma\) is surjective.

Under \(T\mapsto gTh^{-1}\), the right map acts within the full
factor \(\mathbb S_\lambda U\).  Contracting against its dual
therefore leaves exactly the left coefficient line transformed by
\(\mathbb S_\lambda(g)^*\).  Determinant-line scalars and local
lift ambiguities do not change these lines.  This proves the theorem.
\end{proof}

The theorem does not assert that separately projectivized Schur
components always determine their relative scalars.  For webs this
is the recombination problem solved in the separate research archive.  For pencils that issue
disappears altogether.

\subsection{Every pencil, without a regularity hypothesis}

\begin{lemma}[The exterior square of the quadratic representation]
\label{lem:pencil-irreducibility}
For \(n\ge2\), \(\bigwedge^2\Sym^2V=\mathbb S_{(3,1)}V\)
is irreducible.  Consequently
\(\bigwedge^{N-2}\Sym^2V\) is irreducible as well.
\end{lemma}
\begin{proof}
The vector \(x_1^2\wedge x_1x_2\) is a highest-weight vector of
weight \((3,1)\).  By complete reducibility it generates a copy of
\(\mathbb S_{(3,1)}V\).  The hook formula gives
\[
 \dim\mathbb S_{(3,1)}V
   =\frac{n(n-1)(n+1)(n+2)}8
   =\binom{N}{2}.
\]
This exhausts the exterior square.  The last assertion follows by
exterior duality.  These are classical representation identities.
\end{proof}

\begin{theorem}[Intrinsic reconstruction of all quadratic pencils]
\label{thm:natural-pencil-reconstruction}
Let \(n\ge3\) and let \(R,R'\) be two-dimensional subspaces
of \(\Sym^2V\).  For the natural multiplication-failure schemes
of their cube-zero algebras of Hilbert function \((1,n,N-2)\),
\begin{equation}\label{eq:natural-pencil-torelli}
 \widehat D_R\simeq\widehat D_{R'}
 \quad\Longleftrightarrow\quad
 R'=gR\text{ for some }g\in\PGL(V).
\end{equation}
No pencil, base locus, or discriminant is required to be smooth or
regular.  The isomorphism on the left is abstract and carries no
supplied ambient embedding, base, or tensor-factor marking.  All the
reduced schemes are isomorphic to the same Schubert divisor.
\end{theorem}
\begin{proof}
Here \(p=N-2>n\).  By
Theorem~\ref{thm:automatic-quadratic-coefficients} and
Lemma~\ref{lem:pencil-irreducibility}, the normal cone has only one
left coefficient line to recover: the entire line \(\C\beta_R\).
The canonical duality
\begin{equation}\label{eq:pencil-exterior-duality}
 \bigwedge^{N-2}W^*\simeq
       \bigwedge^2W\otimes(\det W)^*
       =\bigwedge^2W\otimes(\det V)^{-(n+1)}
\end{equation}
sends that line to the Pluecker line of \(R\), with the same
projective transformation \(g\).  Indeed the pullback of a
quotient volume is the annihilator-volume dual of the kernel
\(R\).  The common determinant character is a scalar on
projective space.  A nonzero decomposable two-vector recovers its
plane as \(\{w\in W:w\wedge w_R=0\}\).  Thus an abstract
scheme isomorphism implies \(R'=gR\).

Conversely such a linear coordinate change induces the quotient
isomorphism \(S_R\to S_{R'}\), the algebra isomorphism
\(A_R\to A_{R'}\), and the corresponding Grassmannian isomorphism.
It intertwines multiplication and hence preserves its zeroth Fitting
ideal.  This proves the reverse implication.  The assertion about
reductions was proved in \eqref{eq:general-fitting-factorization}.
\end{proof}

\subsection{Discriminants, elementary divisors, and hyperelliptic moduli}

A regular pencil means one containing an invertible symmetric matrix;
it need not have a reduced discriminant.  The classical invariant of
such a pencil consists of its discriminant points together with their
Jordan-block partitions, usually called its Weierstrass--Segre data.
We use the root-labelled partitions, not a pairing of an unlabelled
list with an independently unlabelled divisor.  Their classification
is classical; compare \cite[Theorem~1.1, Corollary~2.1 and \S2]{FevolaMandelshtamSturmfels}.
The following consequence identifies precisely what the intrinsic
failure scheme remembers beyond the discriminant.

\begin{corollary}[Intrinsic recovery of elementary divisors]
\label{cor:intrinsic-segre-data}
For a regular pencil in any dimension \(n\ge3\), the abstract
scheme \(\widehat D_R\) determines all root-labelled elementary
divisors of a matrix pencil, up to projective reparametrization of its
parameter line.  Conversely these data determine its abstract failure
scheme.  In particular the invariant distinguishes different Jordan
partitions over a fixed multiple discriminant root.
\end{corollary}
\begin{proof}
The first implication follows by reconstructing \(R\), then
choosing an invertible member \(A\) and a complementary member
\(B\).  For completeness, the relevant classical sufficiency can
be seen directly.  Put \(C=A^{-1}B\); it is self-adjoint for
\(A\).  If \(hC=C'h\) identifies the Jordan forms for another
pair \((A',B')\), set
\[
 H=A^{-1}h^tA'h.
\]
Then \(H\) is invertible, \(A\)-self-adjoint, and commutes with
\(C\).  There is a polynomial \(k\) in \(H\) with
\(k^2=H\): on each nonzero eigenvalue choose a square root and
truncate its Taylor series to the size of the largest Jordan block,
then combine by the Chinese remainder theorem.  Thus \(k\) is
invertible, self-adjoint, and commutes with \(C\).
The map \(hk^{-1}\) is an isometry from \(A\) to \(A'\)
and intertwines \(C,C'\), so it simultaneously identifies
\(A,B\) with \(A',B'\).

Changing the basis of the pencil replaces \(C\) by a fractional
linear expression with invertible denominator.  At each eigenvalue
its derivative is nonzero; hence each Jordan-block size is preserved
while its root moves by the same projective transformation.  A
parameter change can avoid infinity for both finite spectra.
Root-labelled elementary divisors therefore suffice, and
\eqref{eq:natural-pencil-torelli} proves the converse for the
failure schemes.  This argument recalls the classical classification;
the new input is recovery from an unmarked abstract failure scheme.
\end{proof}

\begin{example}[Identical discriminant, distinct intrinsic failure schemes]
\label{ex:same-discriminant-distinct-failure}
Choose distinct nonzero \(\mu_3,\ldots,\mu_n\), and put
\[
 S_2=\begin{pmatrix}0&1\\1&0\end{pmatrix},\qquad
 J_2=\begin{pmatrix}0&1\\0&0\end{pmatrix}.
\]
Consider the symmetric pairs
\begin{align*}
 A_{\mathrm{s}}&=I_n,&
 B_{\mathrm{s}}&=\diag(0,0,\mu_3,\ldots,\mu_n),\\
 A_{\mathrm{j}}&=S_2\oplus I_{n-2},&
 B_{\mathrm{j}}&=S_2J_2\oplus\diag(\mu_3,\ldots,\mu_n).
\end{align*}
Their determinant binary forms agree up to a nonzero scalar:
\begin{equation}\label{eq:same-discriminant-pencils}
 \det(sA_{\mathrm{s}}+tB_{\mathrm{s}})
 =-\det(sA_{\mathrm{j}}+tB_{\mathrm{j}})
 =s^2\prod_{i=3}^n(s+\mu_i t).
\end{equation}
At the unique double root the ranks are respectively \(n-2\)
and \(n-1\), so no simultaneous projective equivalence exists.
The two root partitions are \((1,1)\) and \((2)\).
Their abstract failure schemes are therefore nonisomorphic although
both their reduced failure schemes and their discriminant divisors,
including multiplicities, agree.  This realizes inside the same
multiplication construction the classical loss of elementary-divisor
information under passage to the determinant.
\end{example}

\begin{corollary}[Hyperelliptic isomorphism classes]
\label{cor:hyperelliptic-failure-invariant}
Fix \(g\ge2\) and \(n=2g+2\).  On the locus of pencils with
simple discriminant, \(\widehat D_R\) gives an injective invariant
of the isomorphism classes of smooth hyperelliptic curves of genus
\(g\).  All these failure schemes have the same reduction; their
nonreduced structures distinguish a \((2g-1)\)-dimensional family
of curve classes.  The pencil invariant also distinguishes the two
regular boundary pencils in
Example~\ref{ex:same-discriminant-distinct-failure}, where their
singular discriminant double covers alone do not.
\end{corollary}
\begin{proof}
The classical curve associated with a simple-discriminant pencil is
the double cover of its parameter line branched at
\(\det(sA+tB)=0\); see \cite[\S4.2, equation (4.3)]{SmithQuadrics}.
Choose an invertible member.  Distinct eigenvalues of \(A^{-1}B\)
have mutually \(A\)-orthogonal eigenlines, each nonisotropic, so
simultaneous diagonalization identifies the pencil with
\(\langle\sum x_i^2,\sum\lambda_i x_i^2\rangle\).
The unordered branch divisor modulo \(\PGL_2\) therefore
determines the pencil, and every reduced divisor of degree \(2g+2\)
arises.  A smooth curve of genus at least two has a unique degree-two
map to a rational curve up to automorphism of that curve: two distinct
such maps would give a birational map to a curve of bidegree \((2,2)\)
in \(\Pj^1\times\Pj^1\), whose arithmetic genus is one.
Thus curve isomorphism is exactly branch-divisor equivalence.
Applying Theorem~\ref{thm:natural-pencil-reconstruction} gives the
asserted injectivity and well-definedness.

Ordered distinct branch points have dimension \(2g+2\).
Normalizing three of them removes the three-dimensional projective
action; taking the finite permutation quotient leaves dimension
\(2g-1\).  Equality of all reductions follows from
\eqref{eq:general-fitting-factorization}.  Finally the two binary
forms in \eqref{eq:same-discriminant-pencils} differ only by a
scalar, a square over \(\C\), and define isomorphic double covers,
whereas their failure schemes have already been separated.
\end{proof}

These are statements about complex isomorphism classes.  They do not
identify moduli stacks, assert a universal family over a coarse moduli
space, or eliminate the different automorphism groups of pencils and
curves.  Nor is the classical pencil--hyperelliptic correspondence a
new Torelli theorem.  The conclusion here is reconstruction of those
objects and of the elementary-divisor information on their boundary
from the abstract nonreduced multiplication-failure scheme, while all
reduced failure data are fixed.
